% IEEE conference paper. Use [journal] for Transactions and other journals.
\documentclass[conference]{IEEEtran}

\usepackage{cite}
\usepackage{amsmath,amssymb}
\usepackage{graphicx}

\begin{document}

\title{An Example Title}

\author{\IEEEauthorblockN{A. U. Thor}
\IEEEauthorblockA{Department, Your University\\
City, Country\\
author@example.org}
\and
\IEEEauthorblockN{Second Author}
\IEEEauthorblockA{Another Institute\\
City, Country\\
second@example.org}}

\maketitle

\begin{abstract}
State the problem, the approach and the main result in one paragraph, with no
citations or symbols that need a definition.
\end{abstract}

\begin{IEEEkeywords}
first keyword, second keyword, third keyword
\end{IEEEkeywords}

\section{Introduction}

Say what is known, what is missing, and what this paper adds. Cite earlier
work by number \cite{einstein1905}, \cite{knuth1984}.

\section{Method}

Describe what you did in enough detail to repeat it.

\begin{equation}
  \bar{x} = \frac{1}{n} \sum_{i=1}^{n} x_i
  \label{eq:mean}
\end{equation}

\begin{table}[htbp]
  \caption{Example Measurements}
  \label{tab:results}
  \centering
  \begin{tabular}{|l|c|c|}
    \hline
    Sample & Mean & Std.\ dev. \\
    \hline
    A & 1.02 & 0.11 \\
    B & 0.97 & 0.08 \\
    \hline
  \end{tabular}
\end{table}

\section{Results}

Report what you found. Table~\ref{tab:results} lists the measurements and
\eqref{eq:mean} defines the mean.

\section*{Acknowledgment}

Thank the people and grants that made the work possible.

\bibliographystyle{IEEEtran}
\bibliography{references}

\end{document}
